\section{Behavior Under Algorithmic Redistricting}
\label{sec:behavior}

\subsection{Why Compactness Produces Near-Zero Partisan Bias}

Partisan Bias is zero for any redistricting plan whose seats-votes function is
symmetric around $(0.50, 0.50)$. The symmetry condition is satisfied when the
distribution of district vote shares is symmetric around 0.50: if the map has
as many districts at $0.50 - \epsilon$ as at $0.50 + \epsilon$ for all $\epsilon$,
then a swing from $v < 0.50$ to $v > 0.50$ that brings the statewide mean from
below to above 50\% will flip the same number of districts from Republican to
Democratic as would be won by Republicans under a symmetric negative swing.

Compactness optimization under \textsc{bisect} produces district vote-share
distributions that are approximately symmetric around the statewide median
Democratic vote share. This occurs because compact, geographically contiguous
districts sample their populations from the surrounding area, and the surrounding
area in most states has a mix of partisan compositions at intermediate distances
from any given center. The result is a distribution of district vote shares that
clusters near the statewide median, with few outliers at the extremes.

The mechanism is not perfect symmetry: NC's urban cores produce outlier districts
at 70--80\% Democratic, creating some right-tail skewness. But the skewness is
much smaller than under enacted maps that deliberately pack Democratic voters into
fewer, higher-margin districts while cracking their remaining voters across many
45--48\% districts. The extreme asymmetry of the enacted distribution ($v_i \in
\{35\%, 38\%, 40\%, 42\%, 44\%, 70\%, 75\%, 80\%\}$, for example) produces large
negative Partisan Bias; the moderate asymmetry of the \textsc{bisect} distribution
produces near-zero Partisan Bias.

\subsection{Symmetry as a Byproduct of Geographic Fairness}

The connection between compactness and partisan symmetry can be stated precisely:
any redistricting plan that allocates population based solely on geographic proximity
(compact, contiguous districts) will tend to have approximately symmetric vote-share
distributions, because geographic proximity in most states does not correlate with
partisan composition in a strongly asymmetric way at the district level.

This means that Partisan Bias near zero is not a design goal of \textsc{bisect}---
the algorithm optimizes population equality and compactness, not partisan symmetry---
but a structural byproduct of those geographic criteria. Expert witnesses can use
this property to make the causal argument: the compactness of \textsc{bisect} maps
drives both their geographic coherence and their partisan near-neutrality. When an
enacted map deviates from compactness in specific ways (districts with narrow necks,
geographic reaching across natural barriers), the compactness deviation and the
partisan asymmetry coincide---because the boundary configurations that compactness
forbids are precisely the configurations that packing and cracking require.

\subsection{Bisect Results}

\textsc{Bisect} maps for NC ($k=14$) produce Bias $\approx -0.01$ ($^\dagger$), which
is below the $1/14 \approx 0.071$ threshold for a one-seat advantage and therefore
rounds to zero in the discrete seat-share framework. WI ($k=8$) and TX ($k=38$)
similarly produce $|\text{Bias}| < 0.02$. The near-zero Bias confirms the compactness
mechanism: geographic boundary-drawing without partisan data produces approximately
symmetric treatment of the two parties at the electoral pivot.

\footnotetext{$^\dagger$ Single-run \textsc{bisect} result using VEST 2020
presidential precinct returns, standard-bisect structure, geographic weights.}
